\chapter{Microphysics: WSM6}\label{ch:mp}

Microphysics turns water vapor into cloud and precipitation and back. WRF offers some thirty schemes. The case uses the
WRF single-moment six-class scheme, WSM6 \citep{hong2006wsm6}, which builds on the revised ice processes of
\citet{hong2004}. Its code lives in \code{phys/physics_mmm/mp_wsm6.F90} (shared
with MPAS), with the WRF wrapper \code{phys/module_mp_wsm6.F} and the driver \code{phys/module_microphysics_driver.F}.

\section{Classes and size distributions}

WSM6 predicts the mixing ratios of six water species: vapor $q_v$, cloud water $q_c$, cloud ice $q_i$, rain $q_r$,
snow $q_s$ and graupel $q_g$. It is \emph{single-moment}: only the mass of each species is predicted. The number of
particles follows from an assumed size distribution. For rain, snow and graupel the distribution is exponential
(Marshall--Palmer),
\begin{equation}
  N_x(D) = N_{0x}\,e^{-\lambda_x D},
\end{equation}
and fixing the intercept $N_{0x}$ makes the slope a function of the mass alone:
\begin{equation}
  \lambda_x = \left(\frac{\pi\,\rho_x\,N_{0x}}{\rho\,q_x}\right)^{1/4}.
\end{equation}
The code writes these slopes as statement functions with two square roots instead of a fourth power (for example
\code{lamdar(x,y) = sqrt(sqrt(pidn0r/(x*y)))}). The parameters, from \code{mp_wsm6.F90} and its initialization
\code{mp_wsm6_init}, are:
\begin{center}
\begin{tabular}{lccc}
\toprule
 & Rain & Snow & Graupel \\
\midrule
Density $\rho_x$ (kg m$^{-3}$) & 1000 & 100 & 500 \\
Intercept $N_{0x}$ (m$^{-4}$) & $8\times10^{6}$ & $2\times10^{6}\,e^{0.12\,(T_0 - T)}$, $\le 10^{11}$ & $4\times10^{6}$ \\
Fall speed $a$, $b$ in $V = a\,D^{b}$ & 841.9, 0.8 & 11.72, 0.41 & 330, 0.8 \\
\bottomrule
\end{tabular}
\end{center}
The snow intercept increases as the temperature falls below $T_0 = 273.15$~K, which mimics smaller, more numerous
crystals in cold air. With the hail option (\code{hail_opt = 1}) the graupel class becomes hail: density 700, intercept
$4\times10^{4}$, $a = 285$.

Cloud water and cloud ice are treated as monodisperse and non-precipitating, apart from a small sedimentation of ice.
The number of ice crystals follows the temperature-independent relation of the scheme,
\begin{equation}
  N_i = \min\bigl(\max\bigl(5.38\times10^{7}\,(\rho\,q_i)^{3/4},\ 10^{3}\bigr),\ 10^{6}\bigr)\ \text{m}^{-3},
\end{equation}
where the code again avoids a power function: it forms $(\rho q_i)^3$ and takes two square roots.

\begin{keypoint}
WSM6 writes its fractional powers as square roots wherever it can. Square root is correctly rounded on CPUs and on
NVIDIA GPUs, so these expressions give identical bits on both without the port's reproducible library. The remaining
powers (about thirty) and exponentials go through \code{rp_pow} and \code{rp_exp} (\cref{ch:repro}).
\end{keypoint}

\section{Processes}

The scheme computes a set of process rates, each a transfer of mass from one species to another, named after the
convention \emph{p} + \emph{target} + \emph{process} + \emph{source}. The ones of WSM6 are:
\begin{center}
\begin{tabular}{ll}
\toprule
Rate & Process \\
\midrule
\code{pcond} & condensation and evaporation of cloud water (saturation adjustment) \\
\code{praut} & autoconversion of cloud water to rain \\
\code{pracw} & accretion of cloud water by rain \\
\code{prevp} & evaporation of rain \\
\code{pigen}, \code{pidep} & initiation of ice crystals; deposition on and sublimation of ice \\
\code{psaut}, \code{pgaut} & autoconversion of ice to snow; of snow to graupel \\
\code{psdep}, \code{pgdep} & deposition on snow; on graupel \\
\code{psacw}, \code{pgacw}, \code{paacw} & accretion of cloud water by snow, by graupel; partition at warm temperatures \\
\code{praci}, \code{psaci}, \code{pgaci} & accretion of cloud ice by rain, snow, graupel \\
\code{piacr}, \code{psacr}, \code{pgacr} & accretion of rain by ice, snow, graupel \\
\code{pracs}, \code{pgacs} & accretion of snow by rain; by graupel \\
\code{psmlt}, \code{pgmlt} & melting of snow; of graupel \\
\code{pseml}, \code{pgeml} & enhanced melting of snow and graupel by accretion of water \\
\code{psevp}, \code{pgevp} & evaporation of melting snow; of melting graupel \\
\bottomrule
\end{tabular}
\end{center}
Comments in the code name the papers each rate comes from. Collection rates integrate the product of two size
distributions and a fall-speed difference. With exponential distributions these integrals reduce to closed forms in
the slopes $\lambda$ and gamma functions of the fall-speed exponents, which the code precomputes in
\code{mp_wsm6_init}.

\section{Sedimentation}

Rain, snow and graupel fall with mass-weighted terminal velocities
\begin{equation}
  V_x = a_x\,\frac{\Gamma(4 + b_x)}{6}\,\lambda_x^{-b_x}\left(\frac{\rho_0}{\rho}\right)^{1/2},
\end{equation}
computed by \code{slope_wsm6}, \code{slope_rain}, \code{slope_snow} and \code{slope_graup}. Falling through thin model
layers at these speeds would require a very short time step with an explicit flux scheme. WSM6 instead uses a
semi-Lagrangian sedimentation (\code{nislfv_rain_plm} for rain, \code{nislfv_rain_plm6} for snow and graupel
together):
\begin{enumerate}
  \item the mass in each layer is moved to where its upper and lower boundaries arrive after one time step;
  \item it is remapped onto the model layers with a piecewise-linear reconstruction;
  \item the mass leaving the lowest layer is the surface precipitation.
\end{enumerate}

\section{One column of WSM6}

\begin{algorithm}[ht]
\caption{One call of \code{mp_wsm6_run} for one column (simplified).}
\label{alg:wsm6}
\begin{algorithmic}[1]
\State convert inputs: $T = \theta\,\pi$ (in the wrapper), density and fall-speed factors
\State $\mathrm{loops} \gets \max(\mathrm{nint}(\Delta t/120\,\mathrm{s}), 1)$ \Comment{1 for both domains of the case}
\For{each minor loop}
  \State saturation mixing ratios over water and ice
  \State slopes and terminal velocities (\code{slope_wsm6})
  \State sedimentation of rain, snow and graupel (\code{nislfv_rain_plm}, \code{nislfv_rain_plm6}); surface precipitation
  \State sedimentation of cloud ice
  \State ice processes below freezing, warm processes above, each limited so no species becomes negative
  \State update $q_x$ and $T$ (latent heating)
  \State saturation adjustment of cloud water (\code{pcond})
\EndFor
\State accumulate surface rain, snow and graupel; convert back $\theta = T/\pi$ (in the wrapper)
\end{algorithmic}
\end{algorithm}

All loops inside \code{mp_wsm6_run} are over $k$ outside and $i$ inside, for a $j$ row of columns. The temporaries are
about eighty arrays dimensioned \code{(its:ite, kts:kte)}, allocated on the stack at every call.

\begin{portnote}
WSM6 is the model's most expensive per-step physics after radiation, and its structure is typical of column physics:
\begin{itemize}
  \item many short loops;
  \item branches on temperature and on species being present;
  \item large temporaries;
  \item a vertical sweep (sedimentation) in each column.
\end{itemize}
The port runs it as one GPU thread per column (\cref{ch:colphys}):
\begin{enumerate}
  \item the wrapper gathers the column's inputs into private arrays;
  \item it calls \code{mp_wsm6_run} with \code{its = ite = 1};
  \item it scatters the results back.
\end{enumerate}
The routine and its callees are compiled for the device, and its temporaries get fixed sizes. The operations per
column are those of the CPU in the same order, so a column harness can prove equality on $10^5$ columns from real
states before the full model is run (T-WSM6-COL, \cref{ch:subsys}). The effective-radius calculation returns
immediately in this configuration (the shortwave scheme does not use it), so the GPU wrapper skips the call.
\end{portnote}

\section{Coupling with the dynamics}

Microphysics runs after the third RK stage on the updated state (\cref{ch:phys}). \code{moist_physics_prep_em}
prepares the physics variables. After the scheme, \code{moist_physics_finish_em} converts the change of temperature
back into a change of $\theta_m$, and stores the diabatic heating rate. The dynamics of the next step uses that rate
(\code{h_diabatic}, in \code{rk_tendency} and \code{small_step_finish}), so the latent heating also acts on the
pressure during the acoustic steps instead of entering as a jump.
